\documentclass[11pt,a4paper]{article}
\usepackage{DDTA_METHODOLOGY_GUIDE_STYLE_R1}

\hypersetup{
  pdftitle={DDTA Documentation Authoring Guide - Revision 7 Rebuild R1},
  pdfauthor={Documentation-Driven Threat Analysis research project}
}

% -----------------------------------------------------------------------------
% PAGE-INTEGRITY HASHES
% Generated by update_page_md5.py over canonical LaTeX payloads between
% PAGE-CONTENT markers. Visible self-hash lines are excluded. The integrity
% index resolves hashes of all listed pages and canonicalizes its own hash as
% <SELF> to avoid circular self-reference.
% -----------------------------------------------------------------------------
\newcommand{\PageMDfive}[1]{\csname PageMDfive#1\endcsname}
%<PAGE-HASH-DEFS>
\expandafter\def\csname PageMDfive1\endcsname{ab8157b85a52435776d7eb2304bfd739}
\expandafter\def\csname PageMDfive2\endcsname{ff9a9668a60997749cf42fa72738d4d5}
\expandafter\def\csname PageMDfive3\endcsname{d731e5773810161fc1612856202f0cf9}
\expandafter\def\csname PageMDfive4\endcsname{3bdac4b5fec97501c23d7968bf1619d5}
\expandafter\def\csname PageMDfive5\endcsname{492447128b2b1b19312002ee80666351}
\expandafter\def\csname PageMDfive6\endcsname{dd6326ef067e8650d8c2fb33662caf1b}
\expandafter\def\csname PageMDfive7\endcsname{8828e008581efc3d40954e198bdb8c0c}
\expandafter\def\csname PageMDfive8\endcsname{690aafdc5bb97b563786dbd7f7e31a59}
%</PAGE-HASH-DEFS>

\newcommand{\PageHashFooter}[1]{%
  \fancyfoot[C]{}%
  \fancyfoot[R]{\sffamily\scriptsize\color{DDTAGray}Pag. \thepage\quad MD5 contenuto: \texttt{#1}}%
}

\newcommand{\IntegrityRow}[3]{#1 & #2 & {\footnotesize\ttfamily #3} \\
}

\begin{document}
\selectlanguage{italian}
\fancyhead[L]{\sffamily\small\color{DDTAGray}DDTA Documentation Authoring Guide}
\fancyhead[R]{\sffamily\small\color{DDTABlue}Revision 7 - Rebuild R1}
\fancyfoot[L]{\sffamily\scriptsize\color{DDTAGray}Documentation authoring methodology}

%<PAGE-DESC:1>Frontespizio della nuova guida di authoring
%<PAGE-CONTENT:1>
\begin{titlepage}
\thispagestyle{empty}
\vspace*{12mm}
{\sffamily\bfseries\color{DDTANavy}\fontsize{15}{18}\selectfont Documentation-Driven Threat Analysis}\par
\vspace{5mm}
{\sffamily\bfseries\color{DDTANavy}\fontsize{28}{32}\selectfont Documentation Authoring Guide}\par
\vspace{2mm}
{\sffamily\color{DDTABlue}\fontsize{18}{22}\selectfont Revision 7 - Rebuild R1}\par

\vspace{15mm}
\begin{tcolorbox}[enhanced,arc=3pt,boxrule=0pt,colback=DDTANavy,left=13pt,right=13pt,top=12pt,bottom=12pt]
{\color{white}\sffamily\bfseries\large Nuova costruzione della guida DDTA per la scrittura della documentazione}\par
\vspace{2mm}
{\color{white!88}\small Il documento viene ricostruito progressivamente da una baseline vuota. Ogni nuova sezione sara' aggiunta solo dopo revisione del significato metodologico e verifica dell'integrita' delle pagine gia' approvate.}
\end{tcolorbox}

\vfill
{\sffamily\scriptsize\color{DDTAGray}MD5 contenuto pagina 1: \texttt{\PageMDfive{1}}}\par
\end{titlepage}
%</PAGE-CONTENT:1>

\clearpage
\setcounter{page}{2}
%<PAGE-DESC:2>Project Problem Framing - definizione e regole di authoring
%<PAGE-CONTENT:2>
\PageHashFooter{\PageMDfive{2}}
\section{Step 0 - Project Problem Framing}

\begin{ddtaopen}[title={RIQUADRO TEMPORANEO - DA CANCELLARE PRIMA DELLA VERSIONE FINALE}]
\textbf{Guida di partenza usata come riferimento genealogico della riscrittura:}
\begin{center}
\texttt{DDTA\_DOCUMENTATION\_AUTHORING\_GUIDE\_R6\_CANDIDATE\_R2.tex}
\end{center}
Questo riferimento serve solo durante la ricostruzione progressiva. Il contenuto della nuova guida non viene importato automaticamente: ogni regola viene riesaminata mentre viene applicata al caso DermaTriage.
\end{ddtaopen}

\subsection{Che cos'e' il framing}
Il \textbf{Project Problem Framing} mette a fuoco il problema --- o la classe di problemi --- che il progetto deve affrontare, insieme al contributo che intende portare entro un boundary sensato. Non e' un tipo documentale L1: resta una \textbf{precondizione dell'authoring DDTA} e non sostituisce MacroRequirement, Decision o FunctionalRequirement.

\begin{ddtaprinciple}[title={Domanda fondamentale}]
Quale problema o classe di problemi deve affrontare il progetto, entro quale boundary significativo, se rimuoviamo temporaneamente le scelte di soluzione correnti?
\end{ddtaprinciple}

\subsection{Regole di authoring}
Quello che va scritto e' il problema, non la pipeline o il modello scelto per risolverlo. Si puo' indicare il risultato atteso a livello di problema, ma senza trasformarlo in una sequenza di requisiti. Le regole operative, per esempio una soglia o un fallback, restano fuori: appartengono a livelli successivi. Ogni significato aggiunto deve essere supportato dalle fonti, e aggiungere dettaglio per riempire i campi non serve.

\begin{ddtacore}[title={Regola di neutralizzazione della soluzione}]
Per riconoscere il problema si mettono \emph{temporaneamente} da parte le scelte tecniche correnti. Questo non autorizza a perdere informazione: le scelte realmente governate verranno recuperate al semantic owner corretto nei livelli successivi.
\end{ddtacore}
%</PAGE-CONTENT:2>

\clearpage
%<PAGE-DESC:3>Project Problem Framing - procedura umana, gate ed esempio DermaTriage
%<PAGE-CONTENT:3>
\PageHashFooter{\PageMDfive{3}}
\subsection{Procedura umana di costruzione}

\begin{enumerate}
  \item \textbf{Leggere le fonti autorizzate prima di scrivere.} Identificare frasi che descrivono scopo, problema, outcome di dominio e boundary, senza assumere che ogni dettaglio sorgente appartenga al framing.
  \item \textbf{Separare problema e soluzione.} Mettere temporaneamente da parte i dettagli di realizzazione, per esempio una pipeline o una soglia, come possibili contenuti di livello successivo.
  \item \textbf{Formulare il problema in linguaggio di dominio.} La frase deve restare vera anche se le scelte tecniche correnti vengono sostituite, salvo che una specifica tecnologia costituisca essa stessa un vincolo governato del problema.
  \item \textbf{Controllare il boundary.} Non ampliare il progetto con obiettivi, attori o responsabilita' che le fonti non autorizzano.
  \item \textbf{Fermarsi quando il problema e' riconoscibile.} MacroRequirement, Decision e FunctionalRequirement arrivano dopo.
\end{enumerate}

\begin{ddtacheck}[title={Gate di accettazione del framing}]
Prima di accettare il testo verificare:
\begin{itemize}
  \item se tutte le scelte tecniche correnti cambiassero, il problema descritto continuerebbe a esistere?
  \item il testo descrive un problema/contributo di progetto oppure sta gia' prescrivendo una soluzione?
  \item sono stati introdotti significati non supportati dalle fonti autorizzate?
  \item ci sono dettagli che possono essere rimandati senza perdere l'identita' del problema?
\end{itemize}
Se una risposta segnala contaminazione o informazione non autorizzata, occorre riformulare prima di procedere ai MacroRequirement.
\end{ddtacheck}

\begin{ddtaanti}[title={Anti-pattern principali}]
\begin{itemize}
  \item trasformare l'architettura corrente nella definizione del problema;
  \item elencare dettagli tecnici solo per rendere il framing piu' concreto;
  \item anticipare fallback o regole condizionali per spiegare casi speciali;
  \item copiare lo stesso significato che verra' poi riscritto integralmente nel primo MacroRequirement;
  \item introdurre obiettivi plausibili di dominio che le fonti non governano.
\end{itemize}
\end{ddtaanti}

\begin{ddtaexample}[title={DermaTriage - framing ottenuto durante l'applicazione della guida}]
\emph{Il problema che DermaTriage deve affrontare e' il supporto al triage precoce di casi dermatologici relativi a lesioni potenzialmente oncologiche. Il progetto parte dalle informazioni gia' disponibili sul caso per determinarne l'urgenza e una priorita' operativa di presa in carico, con l'obiettivo di favorire un instradamento specialistico tempestivo.}
\end{ddtaexample}

Nell'esempio non compaiono dettagli di realizzazione come la pipeline a quattro stadi o il fallback senza immagine. Queste informazioni non vengono eliminate dal progetto: vengono rinviate al livello documentale che ne possiede il significato.
%</PAGE-CONTENT:3>

\clearpage
%<PAGE-DESC:4>Project Problem Framing - LLM Execution Profile e prompt canonico
%<PAGE-CONTENT:4>
\PageHashFooter{\PageMDfive{4}}
\subsection{Profilo di esecuzione LLM}

Il profilo LLM adatta la guida umana all'esecuzione automatica. Non e' una metodologia alternativa e non estende l'authority DDTA. Il suo scopo e' contenere alcune tendenze ricorrenti dei modelli linguistici: completare oltre il necessario, anticipare la soluzione, riempire i campi per inerzia e spacciare una parafrasi per decomposizione.

\subsubsection{Prompt canonico - Project Problem Framing}
\begin{lstlisting}[basicstyle=\ttfamily\footnotesize]
ROLE

You are helping with DDTA documentation authoring. The DDTA Documentation
Authoring Guide is the methodological authority. Work only from the authorized
project sources provided for this step. Do not import project truth from domain
knowledge, earlier reconstructions, Base Analysis, threat models, tooling
requirements, or implementation assumptions.

TASK

Derive the Project Problem Framing.

Framing is a precondition for authoring, not an L1 document type. Identify the
problem or class of problems the project must address, together with the
contribution it intends to make, while setting aside current solution choices.

SOURCE-FIRST RULE

1. Pull out statements supported by the sources that describe purpose, problem,
   domain outcome, actors, and meaningful boundary.
2. Separate problem meaning from current realization details.
3. Keep realization details, but mark them DEFERRED so they can be placed at
   the right semantic level later.

WHAT NOT TO PUT IN THE FINAL FRAMING

- names of components, models, algorithms, or protocols just because they exist;
- pipeline stages or architecture choices;
- thresholds, mappings, fallback rules, or interface details;
- responsibilities or goals the authorized sources do not support;
- details added only to make the framing look complete.

COUNTERFACTUAL CHECK

Ask: "If all current technical and architectural choices were replaced, would
 the stated project problem still exist?" If not, rework the framing unless the
relevant choice is itself a governed problem boundary.

STOP RULE

Stop once the project problem and contribution are recognizable. Do not
pre-write MacroRequirements, Decisions, or FunctionalRequirements.

OUTPUT

Return exactly the analysis schema and final-output schema defined in this
section. Do not omit required fields. For REQUIRED-SLOT fields with no supported
value, use "-".
\end{lstlisting}
%</PAGE-CONTENT:4>

\clearpage
%<PAGE-DESC:5>Project Problem Framing - regola LLM e classificazione dei campi
%<PAGE-CONTENT:5>
\PageHashFooter{\PageMDfive{5}}
\begin{ddtaanti}[title={Regola per l'LLM}]
Un'informazione presente nella source non appartiene automaticamente al framing. Il modello deve prima classificarne il semantic level: conoscere un dettaglio non e' una ragione sufficiente per anticiparlo.
\end{ddtaanti}

\subsection{Schema di output per l'LLM}

Per il framing si tengono separati due livelli: la \textbf{worksheet di analisi} e l'\textbf{output documentale finale}. La worksheet serve a rendere verificabile il ragionamento del modello; non va copiata nella documentazione governata.

\subsubsection{Classificazione dei campi}
\begin{tabularx}{\textwidth}{L{48mm}L{34mm}Y}
\toprule
\textbf{Campo} & \textbf{Classe} & \textbf{Regola} \\
\midrule
Authorized sources & REQUIRED-CONTENT & Elenco delle sole fonti autorizzate per il passo. \\
Source-supported problem meaning & REQUIRED-CONTENT & Fatti di problema direttamente sostenuti dalle fonti. \\
Deferred solution details & REQUIRED-SLOT & Dettagli riconosciuti ma rinviati; ``-'' se assenti. \\
Unsupported / ambiguous claims & REQUIRED-SLOT & Significati non autorizzati o ancora ambigui; ``-'' se assenti. \\
Meaningful boundary & REQUIRED-SLOT & Boundary governato; ``-'' se non stabilito separatamente. \\
Counterfactual result & REQUIRED-CONTENT & PASS o REWORK con motivazione breve. \\
Final framing & REQUIRED-CONTENT & Un blocco conciso di prosa governata. \\
\bottomrule
\end{tabularx}
%</PAGE-CONTENT:5>

\clearpage
%<PAGE-DESC:6>Project Problem Framing - formato esatto dell'output LLM
%<PAGE-CONTENT:6>
\PageHashFooter{\PageMDfive{6}}
\subsubsection{Formato esatto richiesto}
\begin{lstlisting}[basicstyle=\ttfamily\footnotesize]
[DDTA FRAMING ANALYSIS]
AUTHORIZED SOURCES:
- ...

SOURCE-SUPPORTED PROBLEM MEANING:
- ...

MEANINGFUL BOUNDARY:
-

DEFERRED SOLUTION DETAILS:
- ...

UNSUPPORTED / AMBIGUOUS CLAIMS:
- ...

COUNTERFACTUAL CHECK:
PASS | REWORK
Rationale: ...

FRAMING VERDICT:
PROCEED | REWORK | STOP
Rationale: ...

[FINAL PROJECT PROBLEM FRAMING]
<one concise governed prose block>
\end{lstlisting}
%</PAGE-CONTENT:6>

\clearpage
%<PAGE-DESC:7>Project Problem Framing - gate LLM e nota sui controlli quantitativi
%<PAGE-CONTENT:7>
\PageHashFooter{\PageMDfive{7}}
\subsection{Gate sull'output LLM}

\begin{ddtacheck}[title={Controlli prima di accettare l'output LLM}]
\begin{itemize}
  \item il blocco finale contiene solo significato gia' verificato nella worksheet;
  \item i dettagli marcati DEFERRED non ricompaiono nel framing finale senza una ragione esplicita;
  \item nessun campo vuoto viene riempito con una parafrasi di un altro campo;
  \item \texttt{PROCEED} e' ammesso solo se il counterfactual check e la source-authority review sono superati;
  \item l'output LLM resta una proposta da sottoporre alla stessa review semantica prevista per l'autore umano.
\end{itemize}
\end{ddtacheck}

\begin{ddtacore}[title={Nota sui controlli quantitativi}]
Le metriche consolidate di similarita' e retrieval del corpus (Jaccard su shingles, cosine TF--IDF, cosine su embeddings e BM25 per retrieval) verranno definite in una sezione metodologica dedicata. Non sostituiscono i gate semantici e non sono necessarie per accettare il primo framing di un corpus privo di altri elementi documentali.
\end{ddtacore}

\begin{ddtaprinciple}[title={Separazione tra diagnostica e decisione metodologica}]
Le metriche quantitative potranno segnalare sovrapposizioni o possibili rewording, ma non potranno decidere da sole se un elemento DDTA e' corretto. La decisione resta affidata ai gate semantici previsti dalla guida.
\end{ddtaprinciple}
%</PAGE-CONTENT:7>

\clearpage
%<PAGE-DESC:8>Indice di integrita' delle pagine
%<PAGE-CONTENT:8>
\PageHashFooter{\PageMDfive{8}}
\section*{Indice di integrita' delle pagine}
\addcontentsline{toc}{section}{Indice di integrita' delle pagine}

\begin{ddtacore}[title={Regola di verifica}]
Ogni pagina possiede un MD5 del proprio contenuto canonico. Il calcolo usa il contenuto LaTeX compreso tra i marker \texttt{PAGE-CONTENT}; il valore MD5 stampato sulla pagina e' escluso dal payload della pagina stessa. Nell'indice, gli hash delle altre pagine sono risolti prima del calcolo, mentre l'hash della pagina indice viene sostituito dal token canonico \texttt{<SELF>} per evitare una dipendenza circolare.
\end{ddtacore}

\vspace{3mm}
\begin{tabularx}{\textwidth}{L{12mm}L{76mm}Y}
\toprule
\textbf{Pag.} & \textbf{Contenuto} & \textbf{MD5 contenuto} \\
\midrule
%<PAGE-INTEGRITY-ROWS>
\IntegrityRow{1}{Frontespizio della nuova guida di authoring}{\PageMDfive{1}}
\IntegrityRow{2}{Project Problem Framing - definizione e regole di authoring}{\PageMDfive{2}}
\IntegrityRow{3}{Project Problem Framing - procedura umana, gate ed esempio DermaTriage}{\PageMDfive{3}}
\IntegrityRow{4}{Project Problem Framing - LLM Execution Profile e prompt canonico}{\PageMDfive{4}}
\IntegrityRow{5}{Project Problem Framing - regola LLM e classificazione dei campi}{\PageMDfive{5}}
\IntegrityRow{6}{Project Problem Framing - formato esatto dell'output LLM}{\PageMDfive{6}}
\IntegrityRow{7}{Project Problem Framing - gate LLM e nota sui controlli quantitativi}{\PageMDfive{7}}
\IntegrityRow{8}{Indice di integrita' delle pagine}{\PageMDfive{8}}
%</PAGE-INTEGRITY-ROWS>
\bottomrule
\end{tabularx}

\vspace{4mm}
{\small\color{DDTAGray}
A ogni passo, una modifica intenzionale cambia l'MD5 della pagina interessata. Una variazione inattesa dell'MD5 di una pagina gia' approvata impone la revisione della modifica prima di proseguire.}
%</PAGE-CONTENT:8>
\end{document}
